\documentclass[11pt]{article}
\usepackage{amsmath,amssymb,geometry,hyperref}
\geometry{margin=1in}
\title{ROOT V139: From Null-Helix Kinematics to a Testable Gauge Curvature}
\author{Reproducible mathematical audit}
\date{November 23, 2026}
\begin{document}\maketitle
\section{Scope and Units}
Assume $c=1$ only; $\hbar$ and $G$ remain dimensionful. No proof of unification is claimed.
\section{Helix reconstruction}
Let $\mathbf r(t)=(a\cos\omega t,a\sin\omega t,ut)$ and $a^2\omega^2+u^2=1$.
The curvature and torsion are $\kappa=a\omega^2$ and $\tau=u\omega$. Thus
\begin{equation}
\boxed{\omega=\sqrt{\kappa^2+\tau^2},\quad a=\frac{\kappa}{\kappa^2+\tau^2},\quad u=\frac{\tau}{\sqrt{\kappa^2+\tau^2}}.}
\end{equation}
The Darboux angular-velocity vector $\boldsymbol\Omega=\tau\mathbf T+\kappa\mathbf B$ has magnitude $\omega$.
A massless test excitation with $\mathbf p=E\mathbf T$ needs $d\mathbf p/dt=E\kappa\mathbf N$ to follow this non-geodesic curve.
\section{Boost check}
For a boost $\beta$ along the helix axis, define $D=\gamma(1-\beta u)$.
Then $E'=DE$ and $\omega'_{\mathrm{geom}}=\omega_{\mathrm{geom}}/D$; geometric angular frequency does not transform as a plane-wave energy frequency.
\section{Berry curvature}
For $z=(\cos(\theta/2),e^{i\phi}\sin(\theta/2))^T$, set
\begin{equation}
a=-i z^\dagger dz=\frac{1-\cos\theta}{2}d\phi,\qquad
f=da=\frac12\sin\theta\,d\theta\wedge d\phi.
\end{equation}
Equivalently, $f_{\mu\nu}=\frac12\mathbf n\cdot(\partial_\mu\mathbf n\times\partial_\nu\mathbf n)$.
The first Chern number is $(2\pi)^{-1}\int_{S^2}f=1$.
\section{A structural obstruction}
As $f$ is decomposable,
\begin{equation}
\boxed{f\wedge f=0\iff f_{01}f_{23}-f_{02}f_{13}+f_{03}f_{12}=0.}
\end{equation}
A generic Maxwell field may have nonzero $\mathbf E\cdot\mathbf B$, and therefore $F\wedge F\neq0$.
One normalized orientation field cannot supply a universal electromagnetic field strength.
\section{Minimal covariant test action}
\begin{equation}
S=\int d^4x\sqrt{-g}\left[\frac{R}{16\pi G}
-\frac{\hbar}{4g_A^2}F_{\mu\nu}F^{\mu\nu}
-\frac{\hbar f_0^2}{2}(\partial_\mu\theta-A_\mu)(\partial^\mu\theta-A^\mu)\right].
\end{equation}
The coupling $g_A$ is dimensionless; $[f_0]=L^{-1}$, $[A_\mu]=L^{-1}$ and $[\hbar]=ML$.
Stationarity gives
\begin{equation}
\boxed{\nabla_\mu F^{\mu\nu}=-g_A^2f_0^2(\partial^\nu\theta-A^\nu),\qquad
\nabla_\mu(\partial^\mu\theta-A^\mu)=0.}
\end{equation}
For $f_0\neq0$, this is a massive-vector model, not massless electromagnetism.
Pure gauge $A=d\theta$ has $F=0$ and $\partial\theta-A=0$.
\section{Conclusion}
Exact mathematical identities and computational tests hold, but strong, weak, electromagnetic and gravitational interactions have not been derived from one internal helix.
\end{document}
